\documentclass[11pt,a4paper]{article}

\usepackage[utf8]{inputenc}
\usepackage[T1]{fontenc}
\usepackage[margin=2.2cm]{geometry}
\usepackage{parskip}
\usepackage{titlesec}
\usepackage{enumitem}
\usepackage{booktabs}
\usepackage{xcolor}
\usepackage{tcolorbox}
\usepackage{hyperref}
\usepackage{fancyhdr}
\usepackage{listings}
\usepackage{longtable}
\usepackage{array}

\tcbuselibrary{skins,breakable}

\definecolor{accent}{RGB}{30,80,150}
\definecolor{softgray}{RGB}{245,246,248}
\definecolor{codebg}{RGB}{248,248,250}
\definecolor{warn}{RGB}{160,60,40}
\definecolor{good}{RGB}{30,120,70}

\titleformat{\section}{\Large\bfseries\color{accent}}{\thesection}{0.8em}{}
\titleformat{\subsection}{\large\bfseries}{\thesubsection}{0.6em}{}
\titleformat{\subsubsection}{\normalsize\bfseries\itshape}{\thesubsubsection}{0.5em}{}

\hypersetup{
  colorlinks=true, linkcolor=accent, urlcolor=accent,
  pdftitle={AgentIA --- First Measured Baseline}, pdfauthor={AgentIA Engineering}
}

\pagestyle{fancy}
\fancyhf{}
\fancyhead[L]{\small\color{accent}AgentIA --- First Measured Baseline}
\fancyhead[R]{\small\color{accent}\thepage}
\renewcommand{\headrulewidth}{0.4pt}

\lstset{
  basicstyle=\ttfamily\small, backgroundcolor=\color{codebg},
  frame=single, rulecolor=\color{softgray}, breaklines=true,
  columns=fullflexible, keepspaces=true, showstringspaces=false,
  xleftmargin=0.5em, xrightmargin=0.5em,
}

\newtcolorbox{notebox}[1][]{colback=softgray,colframe=accent,boxrule=0.4pt,
  arc=2pt,left=6pt,right=6pt,top=4pt,bottom=4pt,breakable,#1}

\title{\vspace{-1.4cm}\textbf{AgentIA --- First Measured Baseline}\\[0.25cm]
\large The first real generation run: what it cost, how it failed, and what it means
for the optimization approach}
\author{AgentIA Engineering}
\date{2026-09-29}

\begin{document}
\maketitle
\thispagestyle{fancy}

\begin{notebox}
\textbf{What this document records.} On 2026-09-29 the platform executed its first
genuine end-to-end generation session. Before it, no session had ever produced a
measured outcome: the verification metrics, the cost store and the optimization
records were all empty. This is the first data, and the first honest answer to
\emph{how reliable is this agent}.

It also corrects a category error about which gate is which, because that error
would misdirect the optimization design.
\end{notebox}

\begin{notebox}[colframe=warn]
\textbf{CORRECTION, 2026-09-29, after re-measurement. The central finding of this
document is withdrawn.}

It reported \emph{systematic procedural non-compliance}: every stage spending its
entire correction budget, ``not one passed on the first attempt'', and concluded that
this was precisely the signal a reflective skill optimiser needs. That observation
was an artefact of a \textbf{response-parsing bug}, not of model behaviour.

\begin{itemize}[leftmargin=1.4em]
  \item The extractor accepted only a bare JSON object, or code fences carrying the
  artifact path on the fence line. The request never stated a wire format. The shapes
  a real model actually emits --- fenced JSON, prose-wrapped JSON, and code blocks
  whose path sits in a heading --- failed to parse, so the attempt was consumed and
  the stage exhausted \emph{without ever producing a candidate for the gate to judge}.
  \item A second defect recorded \emph{no verdict at all} for a stage that passed on
  its first attempt, so \texttt{passed = None} covered two opposite realities:
  ``nothing was judged'' and ``judged and accepted''. That is why the record read as
  though nothing had been wrong.
\end{itemize}

\textbf{Re-measured after fixing both} --- five blueprints, one session each:
\textbf{25 of 25 stage runs succeeded on their FIRST request}, with zero corrections;
83 artifacts persisted; \textbf{conformance 100 for every session, density 0.0, and
zero findings in every final artifact set}.

\textbf{Consequences.} (i)~There is no procedural non-compliance to correct: the
agent produced compliant output first time, on every stage, for every blueprint.
(ii)~The conformance measure does not vary, so the optimiser's own gate ---
\emph{the measure must vary before anything is built to exploit it} ---
\textbf{fails}, and the optimiser is for that reason not built.
(iii)~The reading in \S\ref{sec:encouraging} and the critical path named in
\S\ref{sec:lost} are withdrawn.

\textbf{What survives, and it is why this document is kept rather than deleted.}
The diagnostic channel did what it was built to do: running the real system found
three defects that a suite of 500+ passing tests could not --- the parser above, a
``verified'' session whose build never ran, and a diagnostic record silently
overwritten by test data. \textbf{The instrument was worth building. The loop it was
built to justify has nothing to optimise.}
\end{notebox}

\tableofcontents
\vspace{0.5cm}

\section{The result}

One session, held-out blueprint \texttt{minimal}, provider \texttt{deepseek},
model \texttt{deepseek-flash}.

\begin{center}
\begin{tabular}{@{}lr@{}}
\toprule
\textbf{Measure} & \textbf{Value} \\
\midrule
Model calls & \textbf{15} \\
Input tokens & 26{,}351 \\
Output tokens & 139{,}679 \\
\textbf{Cost} & \textbf{\$0.085295889} \\
Calls priced & 15 of 15 \\
Terminal state & \textbf{BLOCKED} \\
Artifacts persisted & 0 \\
\bottomrule
\end{tabular}
\end{center}

\textbf{The cost figure is the first the project has ever produced.} Feature 013's
cost tracing works: every call was recorded, priced against the published rate
table, and reconciled. The telemetry, however, did \emph{not} reach MLflow --- see
\S\ref{sec:mlflow}.

\section{How it failed}

Every stage is allowed one initial request plus two corrections. The observed
distribution is the finding:

\begin{center}
\begin{tabular}{@{}lccc l@{}}
\toprule
\textbf{Stage} & \textbf{Calls} & \textbf{Corrections} & \textbf{Outcome} & \textbf{Note} \\
\midrule
SCAFFOLDER & 3 & 2 & passed & 43.4\,s and 11{,}045 output tokens on the rejected first attempt \\
DOMAIN & 3 & 2 & passed & \\
SERVICE & 3 & 2 & passed & \\
CONTROLLER & 3 & 2 & passed & \\
\textbf{TEST} & \textbf{3} & \textbf{2} & \textbf{EXHAUSTED} & \emph{never converged} \\
\bottomrule
\end{tabular}
\end{center}

\textbf{Every stage consumed its entire correction budget. Not one passed on the
first attempt.} This is a systematic, procedural failure pattern, not random
noise --- and it is the single most useful thing the run produced.

\begin{notebox}[colframe=warn]
\textbf{WITHDRAWN.} This claim is false. The stages exhausted because their responses
could not be parsed, not because they were rejected for non-compliance. Re-measured
after the parser was fixed: \textbf{25 of 25 stage runs across five blueprints
succeeded on their first request}, with no corrections at all. There is no
``systematic, procedural failure pattern''; there was a transport bug. See the
correction notice at the front of this document.
\end{notebox}

The terminal error was:

\begin{lstlisting}
[TEST] correction budget exhausted after 2 correction attempts;
human intervention required
\end{lstlisting}

That is constitution Principle V working as specified: bounded autonomy, then
human intervention. The session did not silently succeed, and it did not loop.

\section{The gates are not the same gate}
\label{sec:gates}

A correction that matters for the design. There are \textbf{two unrelated
mechanisms}, both called ``the gate'':

\begin{center}
\begin{tabular}{@{}p{0.20\textwidth}p{0.20\textwidth}p{0.24\textwidth}p{0.24\textwidth}@{}}
\toprule
& \textbf{Compliance gate} & \textbf{Optimization gate} \\
\midrule
Origin & Feature 011 & Feature 014 (SkillOpt) \\
Lives in & \texttt{stages/runner.py} & \texttt{scripts/skillopt/gate.py} \\
Judges & one model \emph{response} & one \emph{skill document} \\
On rejection & forces a retry, bounded at 3 & keeps the current skill \\
Scope & within a session & across sessions \\
Status & \textbf{active --- this is what fired} & deferred, never run \\
\bottomrule
\end{tabular}
\end{center}

\textbf{The gate that rejected fifteen responses is not a SkillOpt-style gate.} It
is the generation-time compliance validator. The SkillOpt gate compares two skill
documents by executing fresh sessions and accepting only a strict improvement; it
has never run.

The distinction is not pedantic. The compliance gate rejects \emph{within} a
session and its verdicts are discarded once the session ends. The optimization
gate rejects \emph{across} sessions and needs its evidence to persist. Confusing
them leads to the belief that the platform already has an optimization loop ---
it does not; it has a validator.

\section{What this means for the reflective approach}

The intended approach is CoEvoSkills-style: a reflective language model proposes
skill edits, guided by structured failure diagnostics, admitted only on measured
evidence. CoEvoSkills reported \textbf{+40.5 points} on SkillsBench with that
shape, and its own ablation is the reason the shape matters:

\begin{center}
\begin{tabular}{@{}lr@{}}
\toprule
\textbf{CoEvoSkills variant} & \textbf{Score} \\
\midrule
Full (co-evolving diagnostic verifier) & 71.1 \\
Opaque pass/fail oracle only & 41.1 \quad ($-30.0$) \\
No evolution at all & 42.4 \quad ($-28.7$) \\
\bottomrule
\end{tabular}
\end{center}

A loop fed only a pass/fail verdict performs no better than no loop.

\subsection{The result is encouraging for the approach}
\label{sec:encouraging}

\begin{notebox}[colframe=warn]
\textbf{WITHDRAWN.} The reasoning below treats the observed failures as real
compliance failures. They were parse failures: no candidate ever reached the gate.
Measured after the fix, the final artifact sets carry \emph{zero} findings, so there
is no recurring rule for a loop to target and no variance for it to move. The gate
fails. This subsection is retained only so the error is not repeated.
\end{notebox}

The observed failure is \emph{exactly the signal such a loop needs}:

\begin{itemize}[leftmargin=1.4em]
  \item The failure is \textbf{recurring} --- it happened on 5 of 5 stages.
  \item It is \textbf{procedural}, not a capability gap: the model can clearly
  produce this code, it simply does not satisfy the house rules on the first or
  second attempt.
  \item It is \textbf{rule-attributed} --- the compliance gate reports which rule
  fired, on which artifact, at which severity.
\end{itemize}

\textbf{The mechanism that is currently rejecting the output is simultaneously the
rich diagnostic channel the optimizer requires.} That is the central opportunity,
and it is why the instrument was built before the optimizer.

\subsection{But the evidence is currently lost}
\label{sec:lost}

\begin{notebox}[colframe=warn]
\textbf{WITHDRAWN, and the diagnosis was wrong in an instructive way.} The artifacts
were \emph{never} lost on the path described here; the records that suggested a loss
were fabricated test output written under the production session key. The real reason
the evidence was empty is that nothing was ever produced to record: the responses
could not be parsed. Naming this the critical path was therefore incorrect, and it is
why the parser bug survived another full baseline. See the correction notice at the
front.
\end{notebox}

The diagnostic record for this session came back
\texttt{artifact\_count = 0}, \texttt{evaluable = false} --- \emph{``nothing to
evaluate''}.

Four stages passed and persisted their artifacts. Their rule violations were
detected, reported, and then \textbf{discarded when the session ended}. A blocked
session therefore produces no evidence at all, which is precisely backwards: the
blocked session is the one most worth diagnosing.

\textbf{Until this is fixed the optimizer has nothing to learn from}, no matter
how good the reflection step is. This is the critical path --- not the loop.

\section{The proposed schema}

\begin{lstlisting}
  session runs
      |
      v
  compliance gate evaluates each stage's candidate
      |                       \
      | (pass)                 \ (fail) -> rule_id, severity, artifact, stage
      v                         v
  artifacts persisted      DIAGNOSTIC CHANNEL          <-- exists (feature 015 US1)
                                |
                                v
                     aggregate by rule across sessions
                     (which rules recur, where, how often)
                                |
                                v
                     reflective LLM proposes skill edits
                     targeting the RECURRING rules
                                |
                                v
                     candidate skill -> fresh sessions -> measure
                                |
                                v
                     admit only on measured contribution;
                     evict skills below the evidence floor
\end{lstlisting}

Mapping onto the existing code:

\begin{center}
\begin{tabular}{@{}p{0.30\textwidth}p{0.34\textwidth}p{0.28\textwidth}@{}}
\toprule
\textbf{Component} & \textbf{Where} & \textbf{State} \\
\midrule
Surrogate verifier (the 30-point component) & \texttt{services/conformance\_diagnostics.py} & \textbf{built} \\
Per-stage attribution & same, reading the stage journal & \textbf{built} \\
Evidence persistence & \texttt{models/diagnostics.py} & built, \textbf{broken on the blocked path} \\
Corpus report & \texttt{services/corpus\_report.py} & built \\
Skill generator (reflection) & \texttt{scripts/skillopt/reflect.py} & exists, \textbf{fed the wrong signal} \\
Evidence gate (leave-one-out) & \texttt{services/skill\_contribution.py} & \textbf{not built} \\
Eviction with a floor & \texttt{scripts/skillopt/currency.py} & \textbf{not built} \\
\bottomrule
\end{tabular}
\end{center}

\section{Next steps, in dependency order}

\begin{enumerate}[leftmargin=1.4em]
  \item \textbf{Capture artifacts on the blocked path.} Without it the diagnostic
  channel is empty exactly when it matters most, and every step below is
  unevaluable. \emph{This is the critical path.}
  \item \textbf{Feed the reflections the diagnostics.} The reflector currently
  receives a derived build exit code, a terminal status and artifact paths. It
  should receive the rule histogram: which rules fired, how often, at which stage.
  The channel exists; the wiring does not.
  \item \textbf{Measure per-skill contribution} by leave-one-out, so an edit can be
  attributed to a rule rather than to a document.
  \item \textbf{Add eviction with an evidence floor.} A library that can only grow
  degrades until injecting a skill is worse than injecting nothing.
  \item \textbf{Grow the task corpus.} Five blueprints cannot support inference: one
  task moves the pass rate 25 points, and the arithmetic floor for any paired
  comparison at that size cannot reach $p<0.05$ under any outcome.
  \item \textbf{Only then} run the optimization loop, scored on conformance rather
  than on build success.
\end{enumerate}

\section{Telemetry: is this in MLflow?}
\label{sec:mlflow}

\textbf{No.} The mirror is wired --- \texttt{cost/mlflow\_sink.py} exposes
\texttt{mirror\_call\_record}, and \texttt{cost/recording.py} calls it on every
recorded call --- but MLflow is \textbf{not installed} in this environment, so the
lazy \texttt{import mlflow} fails and each mirror attempt returns \texttt{False}.

\begin{center}
\begin{tabular}{@{}lll@{}}
\toprule
\textbf{Destination} & \textbf{Status} & \textbf{Role} \\
\midrule
\texttt{backend/cost\_tracking.db} & \textbf{populated} & system of record \\
\texttt{MLFLOW\_TRACKING\_URI} & \texttt{localhost:5000}, unreachable & best-effort mirror \\
\bottomrule
\end{tabular}
\end{center}

This is deliberate: feature 013 made the local store authoritative and the mirror
a non-event when absent, so the cost report is deterministic and works offline. The
figures in this document come from the local store. To see them in MLflow, install
\texttt{mlflow} and run a tracking server; no code change is required.

\section{Defects found by running}

None of the following were reachable from the test suite. All six were found by
executing the real thing, which is the argument for executing it.

\begin{longtable}{@{}p{0.06\textwidth}p{0.50\textwidth}p{0.36\textwidth}@{}}
\toprule
\textbf{\#} & \textbf{Defect} & \textbf{Why the suite missed it} \\
\midrule
\endhead
1 & \textbf{MODEL mode could not run through the graph at all.} LangGraph filters
state to the declared schema, and \texttt{llm\_api\_key} was not declared, so it
was silently stripped between entry and the first node & the 011 tests call
\texttt{run\_stage} with a hand-built state and bypass the graph \\
2 & A table created before a column existed crashed the report:
\texttt{create\_all} creates missing \emph{tables}, never missing \emph{columns} &
the test database is rebuilt fresh every run, so no test saw an older table \\
3 & A crashed session was not recorded, so a batch of four reported three &
tests used well-behaved runners \\
4 & An empty artifact set rendered as \emph{verified clean} & caught by the FR-004
test once written, not before \\
5 & The sandbox volume mount was unreadable on this host (needs \texttt{:Z} under
labelled rootless podman), so Maven found no pom and every session blocked & the
suite mounts nothing real \\
6 & Cost reported \emph{no data} while 15 priced calls sat in the store --- the
session database engine was used to query the cost database & the suite's cost
tests inject the aggregate directly \\
\bottomrule
\end{longtable}

\section{Honest position}

\textbf{What is now known.} The platform generates real code, records real costs,
and refuses to claim success it did not verify. One session costs approximately
\textbf{\$0.085} on the original (broken-parser) run, and \textbf{\$0.025} once the
parser is fixed, because a stage that passes first time costs one request instead of
three. \textbf{The agent does satisfy its own house standard on the first attempt at
every stage} --- the opposite of the claim originally made here.

\begin{notebox}[colframe=warn]
\textbf{CORRECTED.} The sentence that stood here --- ``The agent does not satisfy its
own house standard on the first attempt at any stage, and one stage in five did not
converge'' --- described a parsing failure, not a compliance failure.
\end{notebox}

\textbf{What is not known.} Whether the generated artifacts were any good ---
because they were discarded before they could be diagnosed. The conformance
measure's real-world variance is therefore still unmeasured, and that measurement
is the precondition for the optimization work.

\textbf{What should not be concluded.} That the agent is unreliable in the sense of
being incapable. Fifteen rejections are rejections of \emph{compliance}, not
demonstrations of inability. The distinction is the whole reason procedural
optimization is worth attempting here --- and the whole reason it must be measured
rather than assumed.

\end{document}
